\documentclass{notaaula}

        \titulonota{Magnetismo e força magnética}
        \creditos{Turma 3A · Física · Outubro/2026 · Prof. Josué Cavalcante}

        \begin{document}
        \cabecalho

        \begin{sobre}
        Ímãs, campo magnético e forças sobre cargas e fios percorridos por corrente. Habilidades em foco: EM13CNT304 e EM13CNT308.
        \end{sobre}

        \section{Ímãs e campo magnético}

\begin{copiar}{Polos e linhas de campo}
Todo ímã possui polos norte e sul. Polos iguais se repelem e polos opostos se atraem. Ao cortar um
ímã, surgem ímãs menores com dois polos; não aparecem polos isolados no modelo clássico.

O campo magnético $\vec B$ descreve a ação magnética em cada ponto. Fora de um ímã, as linhas saem
do polo norte e entram no sul; onde estão mais próximas, o campo é mais intenso.
\simbolos{$B$ é medido em tesla (T).}
\diaadia{A bússola se alinha aproximadamente ao campo magnético terrestre.}
\end{copiar}

\section{Força sobre carga em movimento}

\begin{copiar}{Força de Lorentz magnética}
Uma carga $q$ com velocidade $\vec v$ em campo $\vec B$ recebe força magnética de módulo
\[
  F_m=|q|vB\sen\theta.
\]
A força é perpendicular a $\vec v$ e a $\vec B$. Ela é máxima em $90^\circ$ e nula quando a carga
se move paralelamente ao campo. Para carga positiva, use a regra da mão direita; para carga negativa,
inverta o sentido.
\atencao{Como a força é perpendicular à velocidade, um campo magnético isolado não altera a energia cinética da carga.}
\end{copiar}

\begin{exemplo}
Uma carga de módulo $2\times10^{-6}\un{C}$ atravessa perpendicularmente um campo
$B=\dec{0,50}\un{T}$ com $v=3\times10^4\un{m/s}$.
\[
  F_m=2\times10^{-6}\cdot3\times10^4\cdot\dec{0,50}
  =\resultado{3\times10^{-2}\un{N}}.
\]
\end{exemplo}

\section{Trajetória circular}

\begin{copiar}{Campo uniforme}
Se $\vec v\perp\vec B$, a força magnética atua como força centrípeta:
\[
  |q|vB=\frac{mv^2}{r}
  \Rightarrow
  r=\frac{mv}{|q|B}.
\]
Maior massa ou rapidez aumenta o raio; maior carga ou campo diminui o raio. A frequência angular é
$\omega=|q|B/m$ no regime não relativístico.
\end{copiar}

\section{Força sobre corrente}

\begin{copiar}{Fio em campo magnético}
Um trecho retilíneo de comprimento $L$, percorrido por corrente $i$, recebe força
\[
  F=B\,i\,L\sen\theta.
\]
A direção é perpendicular ao fio e ao campo. Esse efeito é a base de motores, alto-falantes e
instrumentos eletromecânicos.
\end{copiar}

\begin{exemplo}[Fio perpendicular]
Um fio de $\dec{0,20}\un{m}$ conduz $4\un{A}$ perpendicularmente a $B=\dec{0,50}\un{T}$:
\[
  F=\dec{0,50}\cdot4\cdot\dec{0,20}
  =\resultado{\dec{0,40}\un{N}}.
\]
\end{exemplo}

\begin{dica}
Faça três passos na regra da mão: identifique primeiro $\vec v$ ou a corrente, depois $\vec B$ e,
por fim, a direção perpendicular. Só então inverta para carga negativa.
\end{dica}

\begin{exercicios}
\nivel{1}{Conceitos}
\begin{questoes}
\item O que ocorre ao cortar um ímã ao meio?
\item Em que situação $F_m=0$ para uma carga em movimento?
\item Em que situação o módulo da força magnética é máximo?
\item O campo magnético sozinho realiza trabalho sobre uma carga? Justifique.
\item Cite duas aplicações da força magnética sobre corrente.
\nivel{2}{Aplicação}
\item Calcule $F_m$ para $q=1\times10^{-6}\un{C}$, $v=2\times10^5\un{m/s}$, $B=\dec{0,30}\un{T}$ e $\theta=90^\circ$.
\item Uma carga entra paralela ao campo. Descreva sua trajetória se nenhuma outra força atuar.
\item Um fio de $\dec{0,50}\un{m}$ conduz $3\un{A}$ perpendicularmente a $B=\dec{0,20}\un{T}$. Ache $F$.
\item No movimento circular, o que acontece com $r$ se $B$ dobra?
\nivel{3}{Síntese}
\item Uma partícula positiva move-se para a direita e $\vec B$ entra na página. Determine o sentido de $\vec F_m$.
\item Compare as trajetórias de duas partículas de mesma carga e rapidez, mas massas $m$ e $2m$.
\item Explique como a força sobre um fio pode produzir rotação em um motor.
\end{questoes}
\gabarito{2) movimento paralelo ou antiparalelo a $\vec B$; 3) $90^\circ$;
6) $6\times10^{-2}\un{N}$; 7) retilínea uniforme; 8) $\dec{0,30}\un{N}$;
9) cai à metade; 10) para cima; 11) a partícula $2m$ descreve raio duas vezes maior.}
\end{exercicios}


\end{document}
